\begin{figure}[h]
\sporkForceStacktrue
\sporkSemanticsStuff
\hrule
\begin{mathpar}
\inferrule*[right=\sporkCSeqStep]{
  \sporkRule{\sporkSub{\sporks_1}{\sporks_2}, \sporkSub{\sporksigma_1}{\sporksigma_2}, \sporkSub{\sporktr^*_1}{\sporktr^*_2}}
}
{
  \sporkRule{\texttt{$\sporkSub{\sporks_1}{\sporks_2}$;$\sporks$}, \sporkSub{\sporksigma_1}{\sporksigma_2}, \sporkSub{\sporktr^*_1}{\sporktr^*_2}}
}
\inferrule*[right=\sporkCSeqDone]{
}
{
  \sporkRule{\sporkSub{\texttt{pass;$\sporks$}}{\sporks}, \sporksigma, \sporktr^*}
}\\
\inferrule*[right=\sporkCIfStep]{
  \sporkRule{\sporkSub{\sporkexpr_1}{\sporkexpr_2}, \sporksigma, \sporkSub{\sporktr^*_1}{\sporktr^*_2}}
}
{
  \sporkRule{\texttt{if $\sporkSub{\sporkexpr_1}{\sporkexpr_2}$:$\sporks$}, \sporksigma, \sporkSub{\sporktr^*_1}{\sporktr^*_2}}
}
\inferrule*[right=\sporkCIfFalse]{
}
{
  \sporkRule{\sporkSub{\texttt{if $\sporkFalse$:$\sporks$}}{\texttt{pass}}, \sporksigma, \sporktr^*}
}\\
\inferrule*[right=\sporkCIfTrue]{
}
{
  \sporkRule{\sporkSub{\texttt{if $\sporkTrue$:$\sporks$}}{\sporks}, \sporksigma, \sporktr^*}
}
\inferrule*[right=\sporkCCudaWarps]{
}
{
  \sporkRule{\sporkSub{\texttt{with CudaWarps(...):$\sporks$}}{\sporks}, \sporksigma, \sporktr^*}
}\\
\inferrule*[right=\sporkCCudaDeviceFunction]{
  % Risk of diverging from the code
  \sporkComment{ (See caption for $\sporkQtlStream$) } \\
  \forall\sporksigma' \in \sporkSigma.\sporktidMap_\mathrm{allThreads}(\sporksigma') = \sporkTid \\
  \sporks_\mathrm{fence1} = \texttt{Fence($\sporkQtlStream \cup \{\texttt{cpu\_in\_order\_qual}\}$, $\sporkQtlStream$)} \# \sporktidMap_\mathrm{allThreads} \\
  \sporks_\mathrm{fence2} = \texttt{Fence($\sporkQtlStream$, $\sporkQtlStream$)} \# \sporktidMap_\mathrm{allThreads}
}
{
  \sporkRule{\sporkSub{\texttt{with CudaDeviceFunction(...):} \sporks_\mathrm{device}}{\texttt{$\sporks_\mathrm{fence1}$;$\sporks_\mathrm{device}$;$\sporks_\mathrm{fence2}$}}, \sporksigma, \sporktr^*}
}\\
\inferrule*[right=\sporkCForStepLo]{
  \sporkRule{\sporkSub{\sporkexpr_1}{\sporkexpr_2}, \sporksigma, \sporkSub{\sporktr^*_1}{\sporktr^*_2}}
}
{
  \sporkRule{\texttt{for $\sporkx$ in $\sporkloopMode$($\sporkSub{\sporkexpr_1}{\sporkexpr_2},\sporkexpr$):$\sporks$}, \sporksigma, \sporkSub{\sporktr^*_1}{\sporktr^*_2}}
}\\
\inferrule*[right=\sporkCForStepHi]{
  \sporkRule{\sporkSub{\sporkexpr_1}{\sporkexpr_2}, \sporksigma, \sporkSub{\sporktr^*_1}{\sporktr^*_2}}
}
{
  \sporkRule{\texttt{for $\sporkx$ in $\sporkloopMode$($\sporkz,\sporkSub{\sporkexpr_1}{\sporkexpr_2}$):$\sporks$}, \sporksigma, \sporkSub{\sporktr^*_1}{\sporktr^*_2}}
}\\
\inferrule*[right=\sporkCForIter]{
  \sporkz_\mathrm{lo} < \sporkz_\mathrm{hi} \\
  \sporkz_\mathrm{next} = \sporkz_\mathrm{lo} + 1 \\
  \sporksigma' = \sporksigma[(\sporkx, ()) \mapsto \sporkz_\mathrm{lo}] \\
  \sporks' = \sporkalphaRename(\sporks)
}
{
  \sporkRule{\sporkSub{\sporkSubNothing}{\texttt{$\sporks'$;}}\texttt{ for $\sporkx$ in $\sporkloopMode$($\sporkSub{\sporkz_\mathrm{lo}}{\sporkz_\mathrm{next}}$,$\sporkz_\mathrm{hi}$):$\sporks$}, \sporkSub{\sporksigma}{\sporksigma'}, \sporktr^*}
}\\
\inferrule*[right=\sporkCForDone]{
  \sporkz_\mathrm{lo} \ge \sporkz_\mathrm{hi}
}
{
  \sporkRule{\sporkSub{\texttt{for $\sporkx$ in $\sporkloopMode$($\sporkz_\mathrm{lo},\sporkz_\mathrm{hi}$):$\sporks$}}{\texttt{pass}}, \sporksigma, \sporktr^*}
}\\
\inferrule*[right=\sporkCAssertStep]{
  \sporkRule{\sporkSub{\sporkexpr_1}{\sporkexpr_2}, \sporksigma, \sporkSub{\sporktr^*_1}{\sporktr^*_2}}
}
{
  \sporkRule{\texttt{assert $\sporkSub{\sporkexpr_1}{\sporkexpr_2}$}, \sporksigma, \sporkSub{\sporktr^*_1}{\sporktr^*_2}}
}
\inferrule*[right=\sporkCAssertPass]{
}
{
  \sporkRule{\sporkSub{\texttt{assert $\sporkTrue$}}{\texttt{pass}}, \sporksigma, \sporktr^*}
}
\end{mathpar}
\caption{Control Semantics.
Note that the loops are stepped sequentially regardless of $\sporkloopMode$,
and that \texttt{with CudaWarps} constructs are equivalent to \texttt{if True:}.
The loop variable is stored in $\sporksigma$ as a 0-dimensional value (read as $\sporkx[]$ by \sporkERead, and visible to $\sporktidMap(\sporksigma)$);
programs never assign loop variables and variable names are assumed unique, so no scope-exit rule is needed.
\texttt{with CudaDeviceFunction} includes implicit ``free'' synchronization between the CPU and all possible CUDA threads.
$\sporkQtlStream$ consists of all ``data RAM'' class $\sporkqtl$ (any proxy),
excluding \texttt{cpu\_in\_order\_qual} and atomic $\sporkqtl$ (Section~\ref{sec:Qtl}).
% Single class in paper corresponds to "mask class" and "proxy class" in code.
No rule exists to step \texttt{assert $\sporkFalse$}.
}
\label{fig:ControlSemantics}
\hrule
\end{figure}
